\documentclass[11pt]{article}
\usepackage{amsmath,amssymb,amsthm}
\usepackage[margin=1.1in]{geometry}
\newtheorem{theorem}{Theorem}
\newtheorem{definition}{Definition}
\newtheorem{proposition}{Proposition}
\theoremstyle{remark}
\newtheorem{remark}{Remark}
\newcommand{\Nbar}{\overline{N}}

\title{The positivity lane:\\ instruments, calibrations, and a requirements
specification\\ for the missing self-adjoint realization}
\author{Carl Gribble\thanks{Programme Note 4, draft v0.1 (WP-RH1--RH3b). Companion to the trilogy
\cite{CompanionI,CompanionII,CompanionIII}, the essay \cite{Essay}, and Programme Notes 1--3
\cite{Defs,FormalLog,Observer}. Developed in collaboration with Claude (Anthropic); every numerical
claim below is reproduced by the cited scripts under \texttt{run\_all.py}, and every scan ran under
the preregistration registry described in Section~2. Responsibility for the content rests with the
author. A full disclosure of AI use appears at the end of the note.}}
\date{August 2026 --- draft}

\begin{document}
\maketitle

\begin{abstract}
In every universe where an analogue of the Riemann Hypothesis is a theorem, the proof runs through a
positive structure: a self-adjoint or unitary realization whose spectrum the zeros are, with
positivity doing the confining --- Castelnuovo positivity for function fields, the adjacency operator
for graphs, the Laplacian for the modular surface. For $\zeta$ that object is missing, and Weil's
criterion makes its absence exact: RH is equivalent to positivity of the explicit-formula functional
$W(g * \tilde g)$ over all admissible test functions. This note documents a measurement programme
aimed at that criterion, run under a preregistration discipline in which gates and predictions are
committed before each instrument's first run and every refuted prediction is permanent in a public
registry. Three instruments and one follow-up were built. (1) The \emph{Weil form on finite windows}:
the Gram matrix of $W$ on a Gaussian family matches its zero-side evaluation to $2.6\times10^{-23}$
and is positive semidefinite at $-8.5\times10^{-24}$; a preregistered prediction about its
near-kernel was refuted by the data, which showed instead that the form on a finite spectral window
is a finite-rank sampler --- resolved rank exactly the number of zeros engaged --- so finite families
are structurally incapable of distinguishing ``RH-positive'' from ``barely positive.''
(2) \emph{Canonical-system reconstruction}: the Jacobi/Krein data of the zero comb is computed
exactly (measure recovered to $1.3\times10^{-38}$), and the arithmetic-free canonical object --- the
closed-form Abel-inverted potential whose WKB counting is the zeros' mean counting function ---
tracks the zeros to RMS $0.54$ with a residual that correlates with a damped prime sum at
$r = 0.9999$: the smooth Hamiltonian misses exactly the primes and nothing else, and the measured
Maslov offset $+0.1251$ matches the predicted $1/8$. (3) The \emph{requirements specification}: an
eight-property, three-universe matrix, every cell machine-checked, adopted with citation, or
measured, with exactly one cell OPEN --- a self-adjoint realization for $\zeta$ --- and two measured
constraints pinned on its occupant: it cannot be smooth, and it must be tight. A follow-up
out-of-sample test earned one further claim: the reconstruction's plateau is the arithmetic-free
two-interval equilibrium of the spectral gap, confirming that the primes live in the corrections.
The note prices its own ambitions honestly: the mechanisms are classical, the tables and the
discipline are the contribution, and no progress on RH itself is claimed.
\end{abstract}

\section{Charter}

Weil's criterion \cite{Weil52,Bombieri} states that RH is equivalent to
\[
W(g * \tilde g) \;\ge\; 0 \qquad \text{for all admissible } g,
\]
where $W$ pairs a test function against the explicit formula's geometric side (archimedean, pole,
and prime terms) and equals, by the Guinand--Weil identity \cite{Guinand,Weil52}, the pairing of
$|\hat g|^2$ against the zero comb. The lane's founding observation is that this criterion has the
same shape as the \emph{proofs} that exist: Weil's function-field proof is positivity of an
intersection form \cite{WeilFF}; the graph-theoretic RH is the spectral bound of the adjacency
operator (Ramanujan $\Leftrightarrow$ Ihara-RH \cite{Ihara,Terras,LPS}); the Selberg zeros are real
because they are eigenvalues of a Laplacian. In each solved case there is a positive structure whose
spectrum the zeros are. For $\zeta$ the structure is missing; the de Branges/Krein theory
\cite{deBranges,Remling} says an equivalent canonical system exists abstractly, and nobody has
exhibited its Hamiltonian.

The lane measures what can be measured about that object, under three standing rules. First, the
words ``proof of RH'' appear in no artifact; the milestones are priced as: instruments and
calibrations, near-certain; novel measured tables and labeled conjectures, likely; a partial
theorem, possible; RH, epsilon. Second, every scan runs under the programme's preregistration
registry: constants, gates, and predictions are committed to the repository \emph{before} first run,
results are append-only, and refuted predictions are permanent. Third, every claim ends in one of
the programme's three states --- ADOPTED (a citation), MACHINE (a reproducible check), or a labeled
conjecture --- and anything surviving a scan climbs the audit ladder before the word ``candidate.''
The Rodgers--Tao theorem $\Lambda_{dB} \ge 0$ \cite{RodgersTao} calibrates the whole enterprise: RH,
if true, is barely true, so no margin-based numerical route can close it; the honest deliverable is
structure, not verification.

\section{Method: the registry}

Each instrument below has a preregistration file and a results file in \texttt{registry/}, tied to
the commit hashes recorded there. The discipline did visible work in this lane, and the note treats
its interventions as content rather than housekeeping:
a committed near-kernel prediction was refuted by the first instrument (Section~3); the second
instrument's first run failed two frozen gates, which diagnosed two bring-up bugs that were fixed
with thresholds untouched and the failures recorded (Section~4); a committed correlation band was
refuted \emph{from above} (Section~4); a constant-mining plan was dissolved by analysis before its
prereg and redirected (Section~6); and a post-hoc reading was refused claim status until it survived
an out-of-sample test on configurations it had never seen (Section~6). All instruments run in
\texttt{run\_all.py}; every number quoted below is printed by the cited script.

\section{The Weil form on finite windows (WP-RH1)}

The instrument (\texttt{weil\_positivity.py}) treats the WP3 pairing engine's identity
\cite{Essay} as a quadratic form: on profiles $\varphi_j(r) = e^{-(r-t_j)^2/4} + e^{-(r+t_j)^2/4}$,
$t_j = 3j$, $j = 0..11$, the Gram matrix $M_{ij} = W(\varphi_i\varphi_j)$ is assembled purely from
the geometric side --- no zeros consumed --- via the closed-form reduction of profile products to the
engine's own test functions. Findings, all under frozen gates:

\emph{Identity and positivity.} All 78 entries match the zero-side evaluation to
$2.6\times10^{-23}$; the minimal eigenvalue is $-8.5\times10^{-24}$, numerical noise: the form is
PSD as far as tested.

\emph{The refutation, and the sampling theorem it exposed.} The committed prediction --- that the
near-kernel is spanned by profiles localized in the zero-free gap $(0, \gamma_1)$ --- was refuted:
the bottom eigenvectors include large-centroid combinations. The measured mechanism: the resolved
rank of $M$ equals the number of zeros the family engages (six, at this family), so the kernel is
\emph{everything that interpolates to zero at the engaged ordinates}. Given the explicit formula
this is immediate ($M$ is a Gram of point evaluations), so it is recorded as known-theorem-derivable
and kept as calibration; its consequence frames the lane: \textbf{positivity on any finite family is
rank-deficient-cheap}, which is a measured restatement of why Weil's criterion quantifies over all
test functions, and of why the information lives in the margin structure as the window grows.

\emph{The prime-free window.} For $g$ supported inside $(-\log 2, \log 2)$ the prime term vanishes
identically and $W = $ archimedean $+$ pole $= -0.37595 + 0.37686 = 9.1\times10^{-4} > 0$, with the
zero side agreeing to $4.2\times10^{-6}$ against a committed tail estimate of $3.8\times10^{-6}$ ---
the unconditionally provable corner of the criterion, exhibited numerically.

\emph{Calibration.} In the universe where the RH-analogue is a theorem (an elliptic curve over
$\mathbb F_{10007}$, Hasse \cite{Hasse}), the Frobenius Toeplitz form came out PSD of \emph{exact
rank 2} at the $10^{-40}$ floor --- positivity with zero margin in all but two directions --- and the
Hasse-violating control fired at $\lambda_{\min} = -38.96$ with growth ratio the golden ratio. The
contrast --- tight finite positivity there, rank-deficient sampling here --- recurs throughout the
lane.

\section{Canonical-system reconstruction (WP-RH2)}

The instrument (\texttt{canonical\_system.py}) computes the two tractable shadows of the missing
Hamiltonian.

\emph{Discrete Krein/Jacobi data.} Stieltjes orthogonalization on the symmetric measure over the
first 40 zeros yields the three-term recurrence coefficients --- the discretized canonical-system
data --- exactly: parity symmetry to $9\times10^{-37}$, and the measure recovered from the
$80\times80$ Jacobi matrix to $1.3\times10^{-38}$. Against two density-matched controls (the smooth
comb $\Nbar^{-1}(k - \tfrac12)$ and a deterministic Poisson comb built from $\pi$'s decimal digits),
the zeros' coefficients hug the smooth comb's \emph{four times tighter} than Poisson's
(fluctuation-ratio $0.246$, committed band $< 0.7$): spectral rigidity, read in Hamiltonian
coordinates. (Adjacent literature: random Jacobi ensembles \cite{KillipNenciu}.)

\emph{The arithmetic-free canonical object.} The Wu--Sprung-type potential \cite{WuSprung,SchHutch}
--- the symmetric well whose WKB counting equals the zeros' mean counting function $\Nbar$ --- is
built by a closed-form Abel inversion,
\[
x(V) \;=\; \frac{1}{\pi}\Bigl[-2u + \sqrt V \,\ln\frac{\sqrt V + u}{\sqrt V - u}\Bigr],
\qquad u = \sqrt{V - 2\pi},
\]
whose defining property (WKB phase $= \Nbar + \tfrac18$ on the built grid) is verified in-code to
$5\times10^{-5}$, and solved forward by Numerov shooting with node-count indexing. The spectrum
counts exactly 60 levels below the committed window and tracks the zeros to RMS $0.54$. The headline
measurement: with $\delta_k = \gamma_k - E_k$, $\bar\rho$ the mean density, and $S$ the damped prime
sum $-\frac1\pi\sum_{n} \frac{\Lambda(n)}{\sqrt n \ln n}\sin(E\ln n)\,e^{-(\ln n)^2/18}$, the
detrended series $\delta_k\bar\rho(\gamma_k)$ and $-S(\gamma_k)$ correlate at
\[
r \;=\; 0.9999,
\]
refuting the committed band $[0.6, 0.95]$ \emph{from above}: $\delta_k\bar\rho \approx \tfrac18 -
S(\gamma_k)$ is numerically an identity (the measured constant $+0.1251$ against the predicted
Maslov/convention offset $\tfrac18$), with no circularity --- the solved spectrum consumes neither
primes nor zeros. In operator language: \textbf{the arithmetic-free Hamiltonian accounts for the
mean density and misses exactly the prime comb, and essentially nothing else.} The mechanism is
Riemann--von Mangoldt plus WKB and belongs to the physics-of-$\zeta$ literature
\cite{WuSprung,SchHutch,BerryKeating}; the quantified table, the controls, and the discipline are
what is claimed.

\emph{Bring-up, disclosed.} The first run failed two frozen gates and thereby found two instrument
bugs: a spurious $1/\pi$ in the preregistered Abel formula (the harmonic-oscillator pair
$x = \sqrt V \leftrightarrow \Phi = E/2$ pins the correct normalization), and a silent
bracket-saturation in the eigenvalue search whose cap masqueraded as a constant residual. Both are
recorded in the results file with thresholds untouched; the headline correlation had already reached
its gate on the broken run through the zeros' own fluctuations, and sharpened by three orders of
tightness once the instrument was sound.

\emph{Termination in the solved universe.} On the elliptic curve's Frobenius measure the same
Stieltjes reconstruction \emph{terminates at step 2 with residual zero}: where the RH-analogue is a
theorem, the canonical system is a finite exact matrix. For $\zeta$ the reconstruction is infinite,
and its truncations' missing content is, by the headline measurement, the primes.

\section{The requirements specification (WP-RH3)}

The instrument (\texttt{requirements\_spec.py}) assembles the property list the missing object must
satisfy as a machine-checked matrix: eight rows (self-adjoint realization; real spectrum; functional
equation; trace formula; mean density law; fluctuations $=$ geodesic/prime comb; rigidity;
tightness) by three columns (elliptic-curve/Frobenius, graph/Ihara, $\zeta$), every cell labeled
ADOPTED, MACHINE, MEASURED, or OPEN. New checks run by this instrument:

\begin{itemize}
\item[--] \emph{Mean density ($\zeta$):} $\max_{k\le100} |\Nbar(\gamma_k) - (k - \tfrac12)| =
0.498$, inside the committed band (Riemann--von Mangoldt \cite{Edwards}).
\item[--] \emph{Rigidity ($\zeta$):} unfolded spacing variance $0.1253$ for the zeros ---
GUE-adjacent (Wigner $\beta=2$: $\approx 0.18$; \cite{Montgomery,Odlyzko}) --- against $1.0755$ for
the $\pi$-digit Poisson control and $0$ for the picket: the ordering the spec requires.
\item[--] \emph{Spectrum $\leftrightarrow$ geodesics (graph, exact):} on the Petersen graph the
integer Hashimoto traces match the Ihara--Bass spectral side \cite{Bass,Terras} to $6\times10^{-29}$
through $m = 12$; and the smooth split $\mathrm{tr}(A^m) = n\,t_m$ (infinite-tree returns
\cite{KestenMcKay}) holds \emph{exactly} below the girth, with first excess $120 = $ the twelve
pentagons traversed ten ways: the graph-universe miniature, as an exact integer identity, of the
previous section's headline --- the first correction to the smooth spectral law counts the shortest
geodesics.
\item[--] \emph{Trace formula (EC, fresh code):} at $p = 13$, the direct count over
$\mathbb F_{169}$ gives $N_2 = 180 = p^2 + 1 - (a^2 - 2p)$.
\end{itemize}

The matrix closes with \textbf{exactly one OPEN cell: (self-adjoint realization, $\zeta$)} --- the
Hilbert--P\'olya object \cite{BerryKeating,Connes} --- and the lane's measurements pin two
constraints on its occupant:
\begin{enumerate}
\item \emph{It cannot be smooth.} An arithmetic-free Hamiltonian reproduces the mean density and
nothing else (Section~4); the prime comb must appear at fluctuation order in the Hamiltonian data.
\item \emph{It must be tight.} Every solved-universe realization is finite and exact with zero
positivity margin to spare (Sections~3--4), and $\zeta$'s margin as measured is $\sim 0$ on every
finite window --- consistent with the ``barely true'' reading \cite{RodgersTao}.
\end{enumerate}

\section{The two-interval plateau, out of sample (WP-RH3b)}

Two disciplinary episodes closed the first-light phase. First, a planned constant-mining target
(PSLQ on the Jacobi tail limits) was dissolved by analysis \emph{before} its prereg: those limits
are the universal two-interval equilibrium of the truncated support
$[-\gamma_N, -\gamma_1] \cup [\gamma_1, \gamma_N]$ --- parity-alternating values
$(\gamma_N \pm \gamma_1)/2$ (classical two-interval recurrence asymptotics; cf.\ \cite{SimonSz}) ---
so mining them would recover $\gamma_1$ and $\gamma_N$ back. Second, WP-RH3's committed window for
testing that identification landed in the truncation-decay regime and its prediction was duly
refuted; the plateau reading visible post hoc was \emph{refused claim status} and converted into an
out-of-sample prereg: the law must predict the plateaus of fresh reconstructions at $N = 50$ and
$N = 60$, targets it had never seen. It did (\texttt{two\_interval\_plateau.py}): deviations
$-1.07/-1.75$ and $-1.43/-2.79$ against a committed band of $\pm 3$, all negative as committed,
odd parity tighter than even as committed. The claim is thereby earned at MACHINE status:
\textbf{the plateau of the reconstructed Hamiltonian is the arithmetic-free two-interval equilibrium
of the spectral gap} --- the leading structure of the canonical-system data depends only on the gap
and the window, so the primes live in the corrections. The systematic negative bias, growing with
$N$ and parity, is recorded as observed-but-unclaimed, a candidate for a density-correction prereg.

\section{Related work, claims, and non-claims}

The mechanisms measured here are classical or belong to active literatures: Weil positivity and its
restricted-class theorems (Connes--Consani \cite{ConnesConsani}); the explicit formula
\cite{Guinand,Weil52,Edwards}; semiclassical potentials for the zeros \cite{WuSprung,SchHutch,
BerryKeating}; canonical systems and de Branges spaces \cite{deBranges,Remling}; GUE statistics of
the zeros \cite{Montgomery,Odlyzko}; graph zeta functions and Ramanujan graphs
\cite{Ihara,Bass,Terras,LPS,Nilli}; crystalline measures at the diffraction frontier
\cite{KurasovSarnak}. This note claims: the assembled instruments and their frozen-gate tables; the
finite-rank sampling framing of finite-window positivity; the rigidity-in-Hamiltonian-coordinates
and two-interval-plateau measured tables (the latter earned out of sample); the requirements matrix
with its named OPEN cell; and the end-to-end demonstration that a preregistration registry can run a
measurement programme against RH-adjacent structure without overclaiming. This note does not claim:
novelty of any mechanism; any progress on RH; and it prices the lane's summit, honestly, at partial
theorems about the OPEN cell's occupant.

\section*{Disclosure of AI use}

In accordance with publisher policies on the use of artificial intelligence, the author discloses
the following. This note and the instruments it reports were developed in an extensive collaboration
with Claude (Anthropic; accessed 2026). The direction of the work --- opening a positivity lane
within the author's research programme, and all decisions about scope, claims, and the
preregistration protocol --- are the author's. Within that direction, the AI system proposed the
lane's design, wrote the preregistrations, instruments, and analysis under the registry discipline
described above (including committing predictions before runs, several of which were refuted and are
so recorded), and drafted this note. Every numerical claim is reproduced by the cited scripts under
\texttt{run\_all.py} in the repository \cite{Code}, and the registry files record the full
prediction-and-refutation history. The author has reviewed the manuscript and takes full
responsibility for its content.

\begin{thebibliography}{99}
\bibitem{Weil52} A. Weil, Sur les ``formules explicites'' de la th\'eorie des nombres premiers,
\emph{Comm. S\'em. Math. Univ. Lund} (1952), 252--265.
\bibitem{Bombieri} E. Bombieri, Remarks on Weil's quadratic functional in the theory of prime
numbers I, \emph{Rend. Mat. Acc. Lincei} \textbf{11} (2000), 183--233.
\bibitem{Guinand} A.~P. Guinand, A summation formula in the theory of prime numbers,
\emph{Proc. London Math. Soc.} \textbf{50} (1948), 107--119.
\bibitem{WeilFF} A. Weil, \emph{Sur les courbes alg\'ebriques et les vari\'et\'es qui s'en
d\'eduisent}, Hermann, Paris, 1948.
\bibitem{Hasse} H. Hasse, Zur Theorie der abstrakten elliptischen Funktionenk\"orper III,
\emph{J. Reine Angew. Math.} \textbf{175} (1936), 193--208.
\bibitem{Ihara} Y. Ihara, On discrete subgroups of the two by two projective linear group over
p-adic fields, \emph{J. Math. Soc. Japan} \textbf{18} (1966), 219--235.
\bibitem{Bass} H. Bass, The Ihara--Selberg zeta function of a tree lattice, \emph{Internat. J.
Math.} \textbf{3} (1992), 717--797.
\bibitem{Terras} A. Terras, \emph{Zeta Functions of Graphs: A Stroll through the Garden}, Cambridge
Univ. Press, 2011.
\bibitem{LPS} A. Lubotzky, R. Phillips, and P. Sarnak, Ramanujan graphs, \emph{Combinatorica}
\textbf{8} (1988), 261--277.
\bibitem{Nilli} A. Nilli, On the second eigenvalue of a graph, \emph{Discrete Math.} \textbf{91}
(1991), 207--210.
\bibitem{KestenMcKay} B.~D. McKay, The expected eigenvalue distribution of a large regular graph,
\emph{Linear Algebra Appl.} \textbf{40} (1981), 203--216.
\bibitem{deBranges} L. de Branges, \emph{Hilbert Spaces of Entire Functions}, Prentice-Hall, 1968.
\bibitem{Remling} C. Remling, \emph{Spectral Theory of Canonical Systems}, de Gruyter, 2018.
\bibitem{WuSprung} H. Wu and D.~W.~L. Sprung, Riemann zeros and a fractal potential, \emph{Phys.
Rev. E} \textbf{48} (1993), 2595--2598.
\bibitem{SchHutch} D. Schumayer and D.~A.~W. Hutchinson, Physics of the Riemann hypothesis,
\emph{Rev. Mod. Phys.} \textbf{83} (2011), 307--330.
\bibitem{BerryKeating} M.~V. Berry and J.~P. Keating, The Riemann zeros and eigenvalue asymptotics,
\emph{SIAM Rev.} \textbf{41} (1999), 236--266.
\bibitem{Connes} A. Connes, Trace formula in noncommutative geometry and the zeros of the Riemann
zeta function, \emph{Selecta Math.} \textbf{5} (1999), 29--106.
\bibitem{ConnesConsani} A. Connes and C. Consani, Weil positivity and trace formula: the
archimedean place, \emph{Selecta Math.} \textbf{27} (2021), 77.
\bibitem{RodgersTao} B. Rodgers and T. Tao, The de Bruijn--Newman constant is non-negative,
\emph{Forum Math. Pi} \textbf{8} (2020), e6.
\bibitem{Montgomery} H.~L. Montgomery, The pair correlation of zeros of the zeta function, in
\emph{Analytic Number Theory}, Proc. Sympos. Pure Math. XXIV, AMS, 1973, 181--193.
\bibitem{Odlyzko} A.~M. Odlyzko, On the distribution of spacings between zeros of the zeta
function, \emph{Math. Comp.} \textbf{48} (1987), 273--308.
\bibitem{KillipNenciu} R. Killip and I. Nenciu, Matrix models for circular ensembles,
\emph{Int. Math. Res. Not.} \textbf{2004}, 2665--2701.
\bibitem{KurasovSarnak} P. Kurasov and P. Sarnak, Stable polynomials and crystalline measures,
\emph{J. Math. Phys.} \textbf{61} (2020), 083501.
\bibitem{SimonSz} B. Simon, \emph{Szeg\H{o}'s Theorem and its Descendants}, Princeton Univ. Press,
2011.
\bibitem{Edwards} H.~M. Edwards, \emph{Riemann's Zeta Function}, Academic Press, 1974.
\bibitem{CompanionI} C. Gribble, Prime counts and prime sums from multiplication tables: a
self-similar formula system, Zenodo (2026), \texttt{doi:10.5281/zenodo.21769103}.
\bibitem{CompanionII} C. Gribble, Certified exact sums of primes over dyadic intervals via a ladder
of generalized factorials, Zenodo (2026), \texttt{doi:10.5281/zenodo.21769105}.
\bibitem{CompanionIII} C. Gribble, Recovering the primes in a dyadic interval from exact moments,
Zenodo (2026), \texttt{doi:10.5281/zenodo.21769107}.
\bibitem{Essay} C. Gribble, The multiplication crystal: dimension, diffraction, and the symmetries
of arithmetic, manuscript (2026), \texttt{papers/multiplication\_crystal\_paper.tex} in \cite{Code}.
\bibitem{Defs} C. Gribble, Crystals and their shadows, Programme Note 1 (2026), in \cite{Code}.
\bibitem{FormalLog} C. Gribble, The multiplication crystal's engine is a formal logarithm,
Programme Note 2 (2026), in \cite{Code}.
\bibitem{Observer} C. Gribble, The observer ladder, Programme Note 3 (2026), in \cite{Code}.
\bibitem{Code} C. Gribble, Reference implementations and preregistration registry, code repository
(2026), \texttt{https://github.com/carlgribble-caa/prime-moments}.
\end{thebibliography}

\end{document}
